\documentclass[10pt,a4paper]{amsart}
\usepackage[margin=19mm]{geometry}
\usepackage{amsmath,amssymb,mathtools}
\usepackage[scr=rsfs,cal=euler]{mathalfa}
\usepackage{xfrac}
\usepackage[unicode,hidelinks,bookmarks=false]{hyperref}
\newcommand{\ie}{i.e.\ }
\newcommand{\cf}{cf.\ }
\newcommand{\resp}{respectively\ }
\newcommand{\Subsection}{\subsection}
\providecommand{\nobreakdash}{\nobreak\hbox{-}}
\setlength{\emergencystretch}{2em}
\multlinegap0pt
\usepackage{amssymb}
\usepackage{enumerate}
\usepackage{mathtools}
\usepackage[shortlabels]{enumitem}
\usepackage{xcolor}
\usepackage{cancel}
\usepackage[normalem]{ulem}
\usepackage{tikz}
\newtheorem{thm}{Theorem}
\title{A strengthening of Chang's lemma}
\author{Gaia Carenini, Leonardo Franchi}
\date{Source: arXiv:2605.07916v1 (CC BY 4.0)}
\begin{document}
\maketitle

\noindent\emph{Source and licence.} A strengthening of Chang's lemma. Version arXiv:2605.07916v1 (CC BY 4.0), distributed by arXiv under the Creative Commons Attribution 4.0 International licence, http://creativecommons.org/licenses/by/4.0/, as recorded in the arXiv licence metadata for this exact version and shown on the abstract page https://arxiv.org/abs/2605.07916v1. Full licence notice: \texttt{licenses/arxiv-2605.07916-CC-BY-4.0.txt}.\par\smallskip

\noindent\textit{Editorial scope.} Counted unit: the article's thm thm:refined-chang-general (source0.tex lines 1063--1091). The article's complete original proof of it is delivered with it (source0.tex lines 1092--1275, one \texttt{proof} environment of 184 lines). No mathematics is written, repaired or re-derived: the statement and the proof are the article's own lines, in the article's own notation, and no retained line is substituted or re-flowed.  No further source-local context is needed: every cross-reference of every retained line resolves inside the two delivered spans.  Nothing was omitted from any retained span, so the package carries no omission record.  No retained line contains a citation, so the wrapper carries no bibliography and no bibliography entry.  No retained line contains a figure environment, an image-inclusion command or drawing code, so no image is delivered and the package declares no asset dependency.  Editorial formatting only, in addition: an AMS \texttt{amsart} preamble that restates the article's own declarations and macros used by the retained lines, namely the environments \texttt{thm} on separate counters, together with the standard packages those lines need.  The retained environments print in this document as Theorem 1 in order of appearance, whereas the article prints the same items with the numbering of its own full document.  The complete native source of the article as distributed in the arXiv e-print for this exact version is bundled unchanged as \texttt{sources/200121/source0.tex}.
\begin{thm}[Refined Chang's lemma for finite abelian groups]
\label{thm:refined-chang-general}
Let $G$ be a finite abelian group, let $A\subseteq G$ be a set of density
$\alpha>0$, and let $\varepsilon\in(0,1)$. Then there exist an integer
$r\ge 0$, a dissociated set $\Delta\subseteq \widehat G$ with $|\Delta|=r$,
and functions $g_\sigma:G\to[0,\infty)$, indexed by $\sigma\in\Sigma_r$,
with
$$
r\le \frac{2}{\varepsilon^2}\log\frac1\alpha.
$$
These functions may be chosen so that, for every $\sigma\in\Sigma_r$,
$$
\mathbb E_{x\in G} g_\sigma(x)=1
\qquad\text{and}\qquad
\operatorname{supp}\widehat{g_\sigma}\subseteq \langle\Delta\rangle,
$$
while, for every $x\in G$,
$
\mathbb E_{\sigma\in\Sigma_r} g_\sigma(x)=1.
$
Moreover, for every $\eta\notin\langle\Delta\rangle$, one has
$$
\mathbb E_{\sigma\in\Sigma_r}
\left|
\mathbb E_{x\sim\mu_A} g_\sigma(x)\,\overline{\eta(x)}
\right|
\le \varepsilon.
$$
\end{thm}

\begin{proof}
We construct inductively dissociated sets
$$
\Delta_0\subseteq \Delta_1\subseteq \cdots
$$
and, for each $i\ge 0$, functions $g_\sigma:G\to[0,\infty)$, indexed by
$\sigma\in\Sigma_i$, such that
$$
\mathbb E_{x\in G}g_\sigma(x)=1,
\qquad
\operatorname{supp}\widehat{g_\sigma}\subseteq \langle\Delta_i\rangle
$$
for every $\sigma\in\Sigma_i$, and
$\mathbb E_{\sigma\in\Sigma_i}g_\sigma(x)=1$ for every $x\in G$. We start with
$\Delta_0=\varnothing$ and $g_\varnothing=1$.

Assume that the construction has reached step $i$. Set
$m_\sigma=\mathbb E_{x\sim\mu_A}g_\sigma(x)$ and define a probability measure
$p_i$ on $\Sigma_i$ by $p_i(\sigma)=2^{-i}m_\sigma$. This is indeed a
probability measure, since
$$
\sum_{\sigma\in\Sigma_i}p_i(\sigma)
=
\mathbb E_{x\sim\mu_A}\mathbb E_{\sigma\in\Sigma_i}g_\sigma(x)
=
1.
$$
Let $u_i$ be the uniform probability measure on $\Sigma_i$, and set
$\Phi_i=D(p_i\|u_i)$. Since, for each fixed $x\in G$, the numbers
$2^{-i}g_\sigma(x)$ form a probability distribution on $\Sigma_i$, and since
$p_i(\sigma)=\sum_x\mu_A(x)2^{-i}g_\sigma(x)$ while
$u_i(\sigma)=\sum_x\mu_G(x)2^{-i}g_\sigma(x)$, the log-sum inequality gives
$$
0\le \Phi_i\le D(\mu_A\|\mu_G)=\log\frac1\alpha.
$$
If, for every $\eta\notin\langle\Delta_i\rangle$, one has
$$
\mathbb E_{\sigma\in\Sigma_i}
\left|
\mathbb E_{x\sim\mu_A}g_\sigma(x)\overline{\eta(x)}
\right|
\le \varepsilon,
$$
then the construction stops. Otherwise choose $\eta\notin\langle\Delta_i\rangle$
such that
$$
\mathbb E_{\sigma\in\Sigma_i}
\left|
\mathbb E_{x\sim\mu_A}g_\sigma(x)\overline{\eta(x)}
\right|
>
\varepsilon.
$$
For each $\sigma\in\Sigma_i$, choose $|\theta_\sigma|=1$ so that
$$
\Re\left(
\theta_\sigma
\mathbb E_{x\sim\mu_A}g_\sigma(x)\overline{\eta(x)}
\right)
=
\left|
\mathbb E_{x\sim\mu_A}g_\sigma(x)\overline{\eta(x)}
\right|,
$$
and put
$
d_\sigma(x)=\Re(\theta_\sigma\overline{\eta(x)}).
$
Equivalently,
$
d_\sigma(x)
=
\frac12\theta_\sigma\overline{\eta(x)}
+
\frac12\overline{\theta_\sigma}\eta(x).
$
Then $|d_\sigma|\le 1$, and, since
$\operatorname{supp}\widehat{g_\sigma}\subseteq\langle\Delta_i\rangle$ while
$\eta,-\eta\notin\langle\Delta_i\rangle$, one has
$$
\mathbb E_{x\in G}g_\sigma(x)d_\sigma(x)
=
\frac12\theta_\sigma\widehat{g_\sigma}(\eta)
+
\frac12\overline{\theta_\sigma}\widehat{g_\sigma}(-\eta)
=
0.
$$
We split each $g_\sigma$ into
$$
g_{\sigma,+}=(1+d_\sigma)g_\sigma,
\qquad
g_{\sigma,-}=(1-d_\sigma)g_\sigma.
$$
These functions are non-negative and have mean $1$. Moreover,
$(g_{\sigma,+}(x)+g_{\sigma,-}(x))/2=g_\sigma(x)$, so, after identifying
$\Sigma_{i+1}$ with $\Sigma_i\times\{\pm1\}$, we still have
$\mathbb E_{\tau\in\Sigma_{i+1}}g_\tau(x)=1$ for every $x\in G$.

Set $\Delta_{i+1}=\Delta_i\cup\{\eta\}$. Since
$\eta\notin\langle\Delta_i\rangle$, the set $\Delta_{i+1}$ is still
dissociated. Also, the identity above shows that the Fourier transform of
$1\pm d_\sigma$ is supported on $\{0,\eta,-\eta\}$. Therefore
$$
\operatorname{supp}\widehat{g_{\sigma,\pm}}
\subseteq
\operatorname{supp}\widehat{g_\sigma}+\{0,\eta,-\eta\}
\subseteq
\langle\Delta_{i+1}\rangle.
$$
Thus the inductive properties are preserved.

It remains to check that the entropy potential increases by a definite amount
at every non-terminal step. Put
$
c_\sigma=
\left|
\mathbb E_{x\sim\mu_A}g_\sigma(x)\overline{\eta(x)}
\right|.
$
Then $0\le c_\sigma\le m_\sigma$, and, by the choice of $\theta_\sigma$,
$\mathbb E_{x\sim\mu_A}g_\sigma(x)d_\sigma(x)=c_\sigma$. Hence
$\mathbb E_{x\sim\mu_A}g_{\sigma,\pm}(x)=m_\sigma\pm c_\sigma$, and therefore
$$
p_{i+1}(\sigma,\pm)
=
\frac12 p_i(\sigma)
\left(1\pm\frac{c_\sigma}{m_\sigma}\right),
$$
with the convention that $c_\sigma/m_\sigma=0$ when $m_\sigma=0$. A direct
calculation gives
$$
\Phi_{i+1}-\Phi_i
=
\mathbb E_{\sigma\sim p_i}
\left[
\frac{1+c_\sigma/m_\sigma}{2}
\log\left(1+\frac{c_\sigma}{m_\sigma}\right)
+
\frac{1-c_\sigma/m_\sigma}{2}
\log\left(1-\frac{c_\sigma}{m_\sigma}\right)
\right].
$$
The function
$$
t\mapsto
\frac{1+t}{2}\log(1+t)
+
\frac{1-t}{2}\log(1-t)
$$
is convex, and at least $t^2/2$ on $[0,1]$. Jensen's
inequality therefore gives
$$
\Phi_{i+1}-\Phi_i
\ge
\frac12
\left(
\mathbb E_{\sigma\sim p_i}\frac{c_\sigma}{m_\sigma}
\right)^2.
$$
Since $p_i(\sigma)=2^{-i}m_\sigma$, we have
$\mathbb E_{\sigma\sim p_i}c_\sigma/m_\sigma
=
\mathbb E_{\sigma\in\Sigma_i}c_\sigma>\varepsilon$. Thus every non-terminal
step gives $\Phi_{i+1}-\Phi_i>\varepsilon^2/2$.

Since $0\le \Phi_i\le \log(1/\alpha)$ at every stage, the construction must
stop after at most
$
\frac{2}{\varepsilon^2}\log\frac1\alpha
$
steps. If $T$ is the first terminal step, set $r=T$ and
$\Delta=\Delta_T$. The terminal condition gives
$$
\mathbb E_{\sigma\in\Sigma_r}
\left|
\mathbb E_{x\sim\mu_A}g_\sigma(x)\overline{\eta(x)}
\right|
\le \varepsilon
\qquad\text{for every }\eta\notin\langle\Delta\rangle,
$$
and all the remaining properties of the functions $g_\sigma$ were preserved
throughout the induction.
\end{proof}

\end{document}
